% !TEX root = ../TCC - CD e IA.tex
\chapter{INTRODUÇÃO}\label{INTRODUÇÃO}

Em um cenário no qual as mídias sociais se consolidaram como plataformas centrais para a comunicação e a formação da opinião pública, o volume massivo de dados gerados pelas interações de usuários \textemdash{} como os comentários no Instagram \textemdash{} emerge como uma fonte inestimável para a análise do engajamento político. Contudo, a natureza não estruturada e a escala desses dados tornam as abordagens analíticas manuais ou isoladas impraticáveis e superficiais \cite{JURAFSKY2025}. A lacuna não está apenas em compreender a polaridade do sentimento, os tópicos subjacentes ou a segmentação do público, mas em converter essa compreensão em decisão prática de comunicação: saber o que produzir, o que amplificar e o que abandonar. Diante disso, esta pesquisa propõe uma abordagem metodológica híbrida \textemdash{} combinando técnicas de Processamento de Linguagem Natural (PLN), como a análise de sentimentos e a modelagem de tópicos, e métodos não supervisionados, como a clusterização \textemdash{} e a articula a um funil de engajamento inspirado no \emph{marketing} de crescimento \cite{FEROZ2024} e nos construtos de \emph{Consumers' Online Brand-Related Activities} (COBRRas) \cite{MUNTINGA2011}. A análise não se limita aos comentários: abrange também as legendas e as transcrições da fala dos vídeos, contrastando o discurso produzido pelas assessorias com a reação expressa pelo público. Aplicada ao estudo de caso dos governadores brasileiros, a proposta visa demonstrar que a integração sinérgica dessas ferramentas supera as limitações das análises tradicionais, entregando não apenas uma leitura mais profunda das dinâmicas de opinião, mas um instrumento replicável de apoio à decisão para o marketing político e a comunicação organizacional.

\section{Contextualização e Definição do Problema}

A ascensão das plataformas de mídias sociais transformou o cenário da comunicação política e da formação da opinião pública. O Instagram, em particular, com sua natureza visual e sua capacidade de gerar engajamento massivo por meio de comentários, curtidas e visualizações, tornou-se um palco onde a percepção da figura política é moldada e expressa. Nesse contexto, a grande quantidade de dados gerados diariamente \textemdash{} textuais e quantitativos \textemdash{} representa um desafio significativo: a extração de inteligência e de \emph{insights} acionáveis de um volume tão grande e heterogêneo de informação.

O problema de pesquisa central reside na incapacidade das metodologias analíticas tradicionais, que frequentemente aplicam técnicas de forma isolada \textemdash{} apenas a análise de sentimento ou apenas a modelagem de tópicos \textemdash{}, de capturar a complexidade e a multidimensionalidade da percepção pública. Tais abordagens fragmentadas falham em identificar as nuances do discurso, os temas específicos que impulsionam um determinado sentimento e, crucialmente, em agrupar publicações e perfis com base em semelhanças que revelem padrões de desempenho e de comportamento comunicacional. Falham, ainda, em confrontar sistematicamente aquilo que a assessoria comunica \textemdash{} nas legendas e nas transcrições de seus vídeos \textemdash{} com aquilo a que o público de fato reage, nos comentários.

Some-se a isso um segundo problema, de ordem prática: mesmo quando esses sinais são medidos, raramente são traduzidos em orientação de decisão. Métricas de superfície, como o volume de curtidas ou de comentários, são frequentemente tomadas como sinônimo de aprovação, ignorando que alto engajamento pode decorrer tanto de aclamação quanto de controvérsia. Assim, a ausência de uma metodologia integrada \textemdash{} que combine sentimentos, tópicos e clusterização em múltiplos níveis de agregação e que organize seus resultados em um funil acionável \textemdash{} deixa um vácuo duplo: no entendimento aprofundado do comportamento do eleitorado e na capacidade de estrategistas de comunicação de tomar decisões informadas em um ambiente digital dinâmico e volátil.

Diante desse vácuo duplo, este trabalho estrutura-se em torno da seguinte pergunta de pesquisa:

\begin{quote}
\emph{Como a integração de análise de sentimentos, modelagem de tópicos e clusterização, aplicada às múltiplas fontes textuais e métricas das publicações de perfis políticos no Instagram, pode ser organizada em um funil de engajamento capaz de orientar as decisões de crescimento (\emph{growth}) de uma assessoria \textemdash{} isto é, o que produzir, o que amplificar e o que abandonar?}
\end{quote}

Dessa questão central desdobram-se três perguntas específicas que guiam o desenvolvimento do estudo. Primeiro, em que medida a abordagem híbrida, ao integrar as três técnicas, revela padrões de percepção pública que as análises isoladas não capturam? Segundo, como o confronto entre o discurso das assessorias \textemdash{} veiculado nas legendas e transcrições \textemdash{} e a reação do público \textemdash{} expressa nos comentários e nas métricas \textemdash{} pode qualificar o diagnóstico de desempenho de um perfil? Terceiro, de que forma os resultados dessa análise podem ser mapeados em estágios de um funil de engajamento, de modo a produzir recomendações práticas e replicáveis para a comunicação política digital?

\section{Relevância e Motivação}

A relevância e a motivação deste estudo residem na lacuna crítica entre a proliferação de dados de mídias sociais e a carência de metodologias analíticas robustas, integradas e \emph{orientadas à decisão} para extrair \emph{insights} significativos. Academicamente, a pesquisa contribui para os campos da Ciência de Dados, da Comunicação Social e do Marketing Político ao propor e validar um \emph{framework} analítico híbrido que transcende as limitações das abordagens tradicionais e isoladas de PLN. Ao demonstrar como a combinação sinérgica de análise de sentimentos, modelagem de tópicos e clusterização pode revelar padrões complexos de opinião \textemdash{} e ao organizar esses padrões segundo os construtos de engajamento dos COBRAs \cite{MUNTINGA2011} e a lógica do funil de \emph{marketing} de crescimento \cite{FEROZ2024} \textemdash{}, o estudo estabelece uma nova fronteira metodológica para a pesquisa aplicada em mídias sociais \cite{PANG2008}.

A motivação principal é suprir uma necessidade prática e urgente: oferecer a estrategistas de marketing político e a comunicadores organizacionais ferramentas que permitam ir além de métricas superficiais, como curtidas e compartilhamentos, proporcionando uma compreensão profunda e, sobretudo, \emph{acionável} da percepção pública e dos temas que efetivamente ressoam com o eleitorado. Em um ambiente político polarizado, no qual a comunicação digital é decisiva, torna-se um diferencial estratégico a capacidade de identificar não apenas o que o público sente, mas por que se sente assim, como os diferentes segmentos de opinião se organizam e \textemdash{} criticamente \textemdash{} em que medida o discurso produzido pela assessoria dialoga com aquilo a que o público de fato reage. O presente trabalho, ao focar nos governadores brasileiros, oferece um estudo de caso oportuno e de grande impacto, ilustrando a aplicação prática da metodologia para informar a tomada de decisão e aprimorar estratégias de engajamento.

\section{Visão Geral do Método}

A metodologia foi concebida como um pipeline de engenharia de dados e mineração de texto, estruturado para processar, de forma integrada, tanto as fontes textuais quanto as métricas quantitativas das publicações. O processo teve início com a coleta, via API da APIFY, de um corpus abrangendo três fontes textuais distintas \textemdash{} os comentários dos usuários, as legendas (\emph{captions}) das publicações e as transcrições da fala dos vídeos \textemdash{} além das métricas de engajamento associadas a cada perfil, publicação do \emph{feed} e \emph{Reel}. A etapa de pré-processamento incluiu a limpeza e a normalização dos textos, com a remoção de \emph{stopwords} da língua portuguesa por meio do recurso correspondente da biblioteca \texttt{nltk}.

Para a análise de sentimentos, empregou-se o \emph{pipeline} de \emph{sentiment-analysis} da biblioteca \texttt{transformers}, com o modelo multilíngue \code{cardiffnlp/twitter-xlm-roberta-base-sentiment}, que classifica cada texto com um rótulo de polaridade (\emph{Positive}, \emph{Negative} ou \emph{Neutral}) e uma pontuação de confiança. A análise foi aplicada de forma diferenciada: aos comentários dos usuários \textemdash{} distinguindo publicações do \emph{feed} e \emph{Reels} \textemdash{} para aferir a reação do público, e às legendas e transcrições para caracterizar o sentimento do discurso produzido pelas assessorias.

A modelagem de tópicos foi realizada com a biblioteca BERTopic, selecionada por sua capacidade de gerar tópicos coerentes e interpretáveis. Para otimizar a representação textual, os \emph{embeddings} de sentenças foram gerados pelo modelo \code{rufimelo/bert-large-portuguese-cased-sts}, pré-treinado em português., pré-treinado em português. A representação dos documentos foi aprimorada com a inclusão de emojis como \emph{tokens}, por meio de sua conversão em texto (\emph{demojize}) e da configuração de um \texttt{CountVectorizer} com expressão regular customizada. Uma inovação da metodologia foi a utilização de um \emph{Large Language Model} (LLM), o Gemini API, para refinar os rótulos dos tópicos: por meio de uma classe \texttt{GeminiDocsRefiner}, uma amostra de comentários de cada tópico foi fornecida ao modelo, que gerou descrições detalhadas, elevando a interpretabilidade dos resultados. A modelagem foi conduzida separadamente sobre os comentários e sobre o discurso oficial (legendas e transcrições), permitindo contrastar os temas que emergem da audiência com aqueles pautados pelas assessorias.

A clusterização foi aplicada em múltiplos níveis de agregação. No nível das publicações, as métricas de engajamento \textemdash{} previamente reduzidas por Análise de Componentes Principais (ACP) \textemdash{} alimentaram um procedimento de clusterização automática, conduzido tanto para os \emph{Reels} quanto para os posts do \emph{feed}, agrupando publicações por semelhança de desempenho. No nível dos perfis, o mesmo princípio foi aplicado às métricas agregadas dos governadores, agrupando-os por padrões de comportamento comunicacional. Por fim, os resultados das análises de sentimento, de tópicos e de clusterização foram integrados e mapeados em um funil de engajamento, articulando os níveis de COBRAs \cite{MUNTINGA2011} aos estágios do funil de crescimento \cite{FEROZ2024}, de modo a converter as evidências analíticas em recomendações práticas.

\section{Contribuições do Trabalho}

As principais contribuições deste trabalho desdobram-se em duas frentes: metodológica e prática. No âmbito metodológico, o estudo inova ao propor, desenvolver e validar uma abordagem analítica híbrida de Processamento de Linguagem Natural (PLN). Diferentemente de pesquisas que aplicam técnicas de forma isolada, esta integra análise de sentimentos, modelagem de tópicos e clusterização em um único pipeline coerente, aplicado a múltiplas fontes textuais \textemdash{} comentários, legendas e transcrições \textemdash{} e a diferentes níveis de agregação \textemdash{} publicações e perfis \textemdash{}, demonstrando que a sinergia dessas ferramentas oferece uma compreensão mais profunda e multifacetada do discurso nas redes sociais. A implementação vale-se de modelos e bibliotecas de ponta, como o \code{cardiffnlp/twitter-xlm-roberta-base-sentiment} para a análise de sentimento e o BERTopic para a modelagem de tópicos. Uma contribuição científica particularmente relevante é a incorporação de um \emph{Large Language Model} (LLM), o Gemini API, para refinar e gerar descrições interpretáveis dos tópicos identificados, elevando a qualidade e a utilidade dos resultados.

No campo prático, a contribuição central é a articulação dessa metodologia a um \emph{framework} acionável de crescimento (\emph{growth}). Ao mapear os resultados analíticos em um funil de engajamento \textemdash{} que integra os níveis de COBRAs \cite{MUNTINGA2011} aos estágios do funil de \emph{marketing} de crescimento \cite{FEROZ2024} \textemdash{}, o trabalho converte descrições da percepção pública em recomendações concretas de comunicação, orientando decisões sobre o que produzir, amplificar e abandonar. O confronto sistemático entre o discurso das assessorias (legendas e transcrições) e a reação do público (comentários e métricas) constitui, nesse sentido, um instrumento de diagnóstico original. Ao aplicar a proposta ao estudo de caso dos governadores brasileiros, a pesquisa fornece um exemplo concreto e replicável de como organizações podem extrair inteligência estratégica de grandes volumes de dados não estruturados do Instagram para aprimorar suas estratégias de engajamento.

\section{Estrutura e Organização do Trabalho}

Este trabalho está organizado em sete capítulos, seguindo uma progressão que parte da fundamentação conceitual, avança pela construção e aplicação da metodologia e culmina na tradução dos resultados em recomendações práticas.

O \textbf{Capítulo 1 \textemdash{} Introdução} contextualiza o problema, enuncia a pergunta de pesquisa e apresenta a relevância, a visão geral do método e as contribuições do estudo. O \textbf{Capítulo 2 \textemdash{} Objetivos} formaliza o objetivo geral e os objetivos específicos que orientam a pesquisa. O \textbf{Capítulo 3 \textemdash{} Referencial Teórico} reúne os fundamentos que sustentam o trabalho, abrangendo o marketing político digital e o \emph{marketing} de crescimento, a teoria de engajamento dos \emph{Consumers' Online Brand-Related Activities} (COBRAs), os precedentes de análise de audiência política em escala e os fundamentos de PLN \textemdash{} análise de sentimentos, modelagem de tópicos e clusterização.

O \textbf{Capítulo 4 \textemdash{} Fontes de Dados} descreve o corpus, sua origem e o processo de coleta das fontes textuais e das métricas de engajamento. O \textbf{Capítulo 5 \textemdash{} Modelagem} detalha o pipeline analítico e apresenta os resultados da análise exploratória, da redução de dimensionalidade, da clusterização, da análise de sentimentos e da modelagem de tópicos. O \textbf{Capítulo 6 \textemdash{} Resultados} consolida a aplicação do \emph{framework} de engajamento, articula os achados ao funil de crescimento e deriva, a partir do estudo de caso dos governadores, recomendações práticas por perfil e por tipo de publicação. Por fim, o \textbf{Capítulo 7 \textemdash{} Conclusões} sintetiza as contribuições, discute as limitações do estudo e aponta direções para trabalhos futuros.

\section{Cronograma e planejamento do trabalho}

A execução deste trabalho foi organizada em fases sequenciais, de modo a garantir a profundidade metodológica exigida por uma abordagem híbrida de mineração de dados. A primeira fase foi dedicada ao planejamento e à estruturação do trabalho, com a definição do problema de pesquisa e a elaboração do referencial teórico que forneceu a base conceitual para a análise.

Em seguida, procedeu-se à coleta e ao pré-processamento dos dados, etapa que envolveu a aquisição, via API, das fontes textuais e das métricas de engajamento dos perfis dos governadores, bem como o tratamento dos dados brutos \textemdash{} limpeza e normalização \textemdash{} para prepará-los para a análise.

O núcleo analítico do projeto correspondeu à fase de análise e modelagem, na qual as técnicas de Processamento de Linguagem Natural (PLN) foram aplicadas de forma integrada: a análise de sentimentos, a modelagem de tópicos e a clusterização, conforme a metodologia proposta. A interpretação dos resultados e a elaboração da discussão permitiram a extração dos \emph{insights} e o mapeamento dos achados no funil de engajamento. Por fim, a última fase foi dedicada à redação, à consolidação das contribuições e à revisão do documento, culminando na defesa do Trabalho de Conclusão de Curso.